
Academic performance has traditionally been associated with the amount of time students dedicate to studying. However, recent perspectives suggest that academic success is influenced by a broader set of factors, including focus, sleep quality, digital habits, and overall well-being \cite{walckiers2017students}. Understanding how these dimensions interact is therefore essential for promoting more efficient and sustainable learning strategies.

To explore the relationship between these variables and academic outcomes, this project will use of the dataset \texttt{student\_records\_missing.csv}, provided by the course instructors, which contains information about students' study habits, lifestyle, and performance indicators. We defined two complementary prediction tasks: a regression task aimed at estimating the \texttt{focus\_index}, and a classification task to predict whether a student is likely to achieve a high \texttt{exam\_score}. A complete machine learning pipeline is implemented, including data preprocessing, baseline model evaluation, advanced model optimization, and clustering analysis.

Additionally, model interpretability techniques are applied to better understand the factors that most influence predictions, and error analysis is conducted to identify limitations and potential improvements. The goal of this work is not only to build accurate predictive models, but also to derive meaningful insights that can support students, educators, and academic institutions in improving learning outcomes.